%!TEX root=./neurips_2023.tex

\section{Related Work}
\label{sec:relatedwork}

\textbf{Model Merging.}
As a method for efficient parameter reuse, model merging is highly practical for combining model skills without the need for additional training. Its applications include merging checkpoints of models during training to enhance model performance~\citep{Gupta2020Stochastic,wortsman2022model}, integrating models with varying hyperparameter settings to improve generalization~\citep{arpit2022ensemble,rame2022diverse}, and merging pre-trained model with fine-tuned model to enhance robustness~\citep{wortsman2022robust}.
 
The focus of research on the model merging algorithm lies in the coefficient weights assigned to each parameter in the model.
\citet{matena2022merging} use Fisher matrix to measure the difference in information contained in each parameter of models. 
RegMean\citep{jin2022dataless} are designed to ensure the merged model performs like model A on data A and model B on data B. 
Although both methods outperform the simple average method, the complexity involved in computing and transferring on data means that research on scenarios for merging models generally still utilizes the simple average method.

The previously described method for model merging only utilized the addition operation. 
\citet{ilharco2022editing} proposed a method where new vectors are constructed for different tasks by adding and subtracting task vectors. The task vectors are created by subtracting the fine-tuned model parameters from the pre-trained model parameters. Therefore, the task vectors contain the same quantity of parameters as the large-scale pre-trained language model. Nonetheless, we achieved comparable results using lightweight PEMs.

\textbf{Modular Combining. }
In recent years, multitask learning which combines modular modules at inference time has gained popularity. 
Various types of modular modules have been introduced, such as Adapter ~\citep{wang2022adamix}, Demix layer~\citep{gururangan-etal-2022-demix} that replaces the traditional FFN layer in transformers as a task vector, and skill vector~\citep{ponti2022combining} which refers to skill instead of task. Nevertheless, training access to all data has weakened its applicability.
AdapterFusion utilize the off-the-shelf adapter without any further training ~\citep{pfeiffer2021adapterfusion}. However, the fusion head still requires few-shot training, unlike our zero-shot setting, which simplify the process by employing direct usage without any need for extra data and training.

\citet{qin2022exploring} explored the performance of Pfeiffer adapter merging on multiple hyperparameters and multitask scenarios; they reported a light performance drop compared to full fine-tuning merging.
\citet{chronopoulou2023adaptersoup} have achieved good results in language modeling by using the Houlsby adapter multi-domain merging.
Unlike the two prior studies that solely employed addition operations, we merge LoRA~\citep{hu2021lora} and \ia~\citep{liu2022few}, define both  addition and subtraction, and perform multiple experiments on complex combination methods across different types of tasks. 
The results indicate that the approaches are comparable to fully fine-tuned merged models and extend the compatibility of parameter-efficient modules (PEM) from BERT-like models to LLM such as Llama.
